% 自动生成：问题四全部图表说明。
% 本文件不是正文强制内容，而是用于从大量备选图中挑选论文插图。
\section*{问题四图表说明与选图建议}
本次程序实际成功生成\textbf{65}幅不同图表。每幅图均给出图意、可支持的结论以及是否建议放入正文；未被程序实际生成的可选地理图不会列入本文件。
图表数量以\texttt{Q4\_图表清单.csv}和磁盘中实际存在的文件为准，不再使用预先写死的数量。

\subsection*{图 1：问题四总体求解流程}
\textbf{文件：}\texttt{fig\_q4\_01\_solution\_flow.png}；矢量版：\texttt{fig\_q4\_01\_solution\_flow.pdf}\\
\textbf{图中内容：}从问题三最终输出出发，依次完成任务压缩、超图建模、严格冲突判定、完整枚举和区间资源配置。\\
\textbf{主要意义：}概括问题四与问题三的数据继承关系，并说明本题没有重新优化运输路线或通信位置。\\
\textbf{正文建议：}正文强烈推荐

\subsection*{图 2：严格继承与按组复制两种情景}
\textbf{文件：}\texttt{fig\_q4\_02\_scenario\_definition.png}；矢量版：\texttt{fig\_q4\_02\_scenario\_definition.pdf}\\
\textbf{图中内容：}明确区分题面字面严格继承和用于完成2组/3组资源核算的中继任务完整复制扩展口径。\\
\textbf{主要意义：}避免把结构性不可行与库存不足混为一谈，也避免用增加设备数量掩盖原中继架次跨组共享问题。\\
\textbf{正文建议：}正文强烈推荐

\subsection*{图 3：各运输不可拆分单元包含的服务区数量}
\textbf{文件：}\texttt{fig\_q4\_03\_unit\_service\_count.png}；矢量版：\texttt{fig\_q4\_03\_unit\_service\_count.pdf}\\
\textbf{图中内容：}按同一运输架次多服务区同组约束压缩后，统计每个单元包含的服务区数量。\\
\textbf{主要意义：}验证15个服务区被压缩为9个不可拆分单元，并识别由多点运输架次形成的成对单元。\\
\textbf{正文建议：}正文备选

\subsection*{图 4：各运输不可拆分单元承担的运输架次数}
\textbf{文件：}\texttt{fig\_q4\_04\_unit\_trip\_count.png}；矢量版：\texttt{fig\_q4\_04\_unit\_trip\_count.pdf}\\
\textbf{图中内容：}统计问题三25个运输架次在9个不可拆分单元上的分布。\\
\textbf{主要意义：}服务区数量相同的单元可能承担完全不同的运输频次，说明分区均衡不能只看服务区个数。\\
\textbf{正文建议：}备选

\subsection*{图 5：各运输不可拆分单元的运输作业工作量}
\textbf{文件：}\texttt{fig\_q4\_05\_unit\_transport\_workload.png}；矢量版：\texttt{fig\_q4\_05\_unit\_transport\_workload.pdf}\\
\textbf{图中内容：}将属于同一单元的固定运输架次持续时间累加。\\
\textbf{主要意义：}为组间工作量均衡指标提供最小粒度的解释基础。\\
\textbf{正文建议：}正文备选

\subsection*{图 6：中继架次—运输单元关联矩阵}
\textbf{文件：}\texttt{fig\_q4\_06\_hypergraph\_incidence.png}；矢量版：\texttt{fig\_q4\_06\_hypergraph\_incidence.pdf}\\
\textbf{图中内容：}以13个中继架次为行、9个运输不可拆分单元为列，标出实际Q3通信保障关系。\\
\textbf{主要意义：}矩阵中的多列非零行就是超边；它揭示共享中继把多个运输单元绑定在一起，是严格分区冲突的直接证据。\\
\textbf{正文建议：}正文强烈推荐

\subsection*{图 7：运输单元—中继架次二部关联图}
\textbf{文件：}\texttt{fig\_q4\_07\_unit\_relay\_bipartite.png}；矢量版：\texttt{fig\_q4\_07\_unit\_relay\_bipartite.pdf}\\
\textbf{图中内容：}将运输不可拆分单元与问题三实际使用的中继架次置于两侧，连线表示该中继实际保障了该单元中的运输任务。\\
\textbf{主要意义：}比普通两两网络更直观地展示高阶共享关系；可观察到关联链条贯通全部单元，因此严格继承时不能切成2组或3组。\\
\textbf{正文建议：}正文强烈推荐

\subsection*{图 8：严格继承口径下的运输单元连通图}
\textbf{文件：}\texttt{fig\_q4\_08\_strict\_connectivity.png}；矢量版：\texttt{fig\_q4\_08\_strict\_connectivity.pdf}\\
\textbf{图中内容：}把同一实际中继架次连接的运输单元两两展开，仅用于显示连通性。\\
\textbf{主要意义：}9个运输单元处于同一个连通分量，说明若原中继架次不可复制且资源不得跨组，则K=2、3均结构性不可行。\\
\textbf{正文建议：}正文强烈推荐

\subsection*{图 9：完整枚举的无标签分区数量}
\textbf{文件：}\texttt{fig\_q4\_09\_partition\_counts.png}；矢量版：\texttt{fig\_q4\_09\_partition\_counts.pdf}\\
\textbf{图中内容：}展示9个不可拆分单元分别划分为2组和3组的无标签集合分区数量。\\
\textbf{主要意义：}K=2仅255个、K=3仅3025个，规模很小，因此完整枚举比使用随机启发式更可靠、可复现。\\
\textbf{正文建议：}正文推荐

\subsection*{图 10：各运输单元关联的实际中继架次数}
\textbf{文件：}\texttt{fig\_q4\_10\_unit\_relay\_degree.png}；矢量版：\texttt{fig\_q4\_10\_unit\_relay\_degree.pdf}\\
\textbf{图中内容：}统计每个运输单元实际使用过多少个问题三中继架次。\\
\textbf{主要意义：}高关联度单元更容易在分区时引起中继任务复制，是通信资源压力的重要来源。\\
\textbf{正文建议：}备选

\subsection*{图 11：各中继架次连接的运输单元数量}
\textbf{文件：}\texttt{fig\_q4\_11\_relay\_hyperedge\_size.png}；矢量版：\texttt{fig\_q4\_11\_relay\_hyperedge\_size.pdf}\\
\textbf{图中内容：}以实际通信保障关系统计每个中继架次同时关联多少个运输不可拆分单元。\\
\textbf{主要意义：}规模大于1的超边是潜在跨组复制源；规模为1的中继任务不会因该超边自身导致额外复制。\\
\textbf{正文建议：}正文备选

\subsection*{图 12：K=2 全部分区的中继无人机总需求分布}
\textbf{文件：}\texttt{fig\_q4\_12\_K2\_relay\_drone\_hist.png}；矢量版：\texttt{fig\_q4\_12\_K2\_relay\_drone\_hist.pdf}\\
\textbf{图中内容：}对K=2的全部255个无标签分区统计独立配置所需中继无人机总数量。\\
\textbf{主要意义：}展示中继资源需求不是单一数值，而是由分区方式决定的离散分布；主方案可置于整体分布中理解。\\
\textbf{正文建议：}正文备选

\subsection*{图 13：K=3 全部分区的中继无人机总需求分布}
\textbf{文件：}\texttt{fig\_q4\_13\_K3\_relay\_drone\_hist.png}；矢量版：\texttt{fig\_q4\_13\_K3\_relay\_drone\_hist.pdf}\\
\textbf{图中内容：}对K=3的全部3025个无标签分区统计独立配置所需中继无人机总数量。\\
\textbf{主要意义：}展示中继资源需求不是单一数值，而是由分区方式决定的离散分布；主方案可置于整体分布中理解。\\
\textbf{正文建议：}正文备选

\subsection*{图 14：K=2 全部分区的中继能源组件需求分布}
\textbf{文件：}\texttt{fig\_q4\_14\_K2\_relay\_energy\_hist.png}；矢量版：\texttt{fig\_q4\_14\_K2\_relay\_energy\_hist.pdf}\\
\textbf{图中内容：}对K=2全部候选统计中继能源组件独立配置规模。\\
\textbf{主要意义：}显示能源组件缺口与中继无人机缺口并非完全同步，应分类核算。\\
\textbf{正文建议：}备选

\subsection*{图 15：K=3 全部分区的中继能源组件需求分布}
\textbf{文件：}\texttt{fig\_q4\_15\_K3\_relay\_energy\_hist.png}；矢量版：\texttt{fig\_q4\_15\_K3\_relay\_energy\_hist.pdf}\\
\textbf{图中内容：}对K=3全部候选统计中继能源组件独立配置规模。\\
\textbf{主要意义：}显示能源组件缺口与中继无人机缺口并非完全同步，应分类核算。\\
\textbf{正文建议：}备选

\subsection*{图 16：K=2 中继任务新增副本数分布}
\textbf{文件：}\texttt{fig\_q4\_16\_K2\_copy\_hist.png}；矢量版：\texttt{fig\_q4\_16\_K2\_copy\_hist.pdf}\\
\textbf{图中内容：}按超图connectivity-1口径统计每个分区需要增加的中继任务副本数。\\
\textbf{主要意义：}刻画分区对原中继共享关系的破坏程度，是比单纯设备数量更直接的结构代价。\\
\textbf{正文建议：}正文推荐

\subsection*{图 17：K=3 中继任务新增副本数分布}
\textbf{文件：}\texttt{fig\_q4\_17\_K3\_copy\_hist.png}；矢量版：\texttt{fig\_q4\_17\_K3\_copy\_hist.pdf}\\
\textbf{图中内容：}按超图connectivity-1口径统计每个分区需要增加的中继任务副本数。\\
\textbf{主要意义：}刻画分区对原中继共享关系的破坏程度，是比单纯设备数量更直接的结构代价。\\
\textbf{正文建议：}正文推荐

\subsection*{图 18：K=2 中继复制代价—运输工作量均衡关系}
\textbf{文件：}\texttt{fig\_q4\_18\_K2\_copy\_balance\_scatter.png}；矢量版：\texttt{fig\_q4\_18\_K2\_copy\_balance\_scatter.pdf}\\
\textbf{图中内容：}每个点对应一个完整分区；横轴为中继复制架次数，纵轴为固定运输工作量CV。\\
\textbf{主要意义：}直观显示减少跨组中继切割与提高工作量均衡之间的真实权衡。\\
\textbf{正文建议：}正文强烈推荐

\subsection*{图 19：K=3 中继复制代价—运输工作量均衡关系}
\textbf{文件：}\texttt{fig\_q4\_19\_K3\_copy\_balance\_scatter.png}；矢量版：\texttt{fig\_q4\_19\_K3\_copy\_balance\_scatter.pdf}\\
\textbf{图中内容：}每个点对应一个完整分区；横轴为中继复制架次数，纵轴为固定运输工作量CV。\\
\textbf{主要意义：}直观显示减少跨组中继切割与提高工作量均衡之间的真实权衡。\\
\textbf{正文建议：}正文强烈推荐

\subsection*{图 20：K=2 中继无人机配置—运输工作量均衡关系}
\textbf{文件：}\texttt{fig\_q4\_20\_K2\_relay\_balance\_scatter.png}；矢量版：\texttt{fig\_q4\_20\_K2\_relay\_balance\_scatter.pdf}\\
\textbf{图中内容：}将实体中继资源最少配置与运输工作量均衡指标直接关联。\\
\textbf{主要意义：}说明理论最少中继方案不一定是主提交方案；主方案还需兼顾运输侧零增配与均衡。\\
\textbf{正文建议：}正文推荐

\subsection*{图 21：K=3 中继无人机配置—运输工作量均衡关系}
\textbf{文件：}\texttt{fig\_q4\_21\_K3\_relay\_balance\_scatter.png}；矢量版：\texttt{fig\_q4\_21\_K3\_relay\_balance\_scatter.pdf}\\
\textbf{图中内容：}将实体中继资源最少配置与运输工作量均衡指标直接关联。\\
\textbf{主要意义：}说明理论最少中继方案不一定是主提交方案；主方案还需兼顾运输侧零增配与均衡。\\
\textbf{正文建议：}正文推荐

\subsection*{图 22：K=2 总资源缺口—运输工作量均衡关系}
\textbf{文件：}\texttt{fig\_q4\_22\_K2\_shortage\_balance\_scatter.png}；矢量版：\texttt{fig\_q4\_22\_K2\_shortage\_balance\_scatter.pdf}\\
\textbf{图中内容：}按资源类别分别计算缺口后求和，仅用于散点横轴，不把不同资源成本等价化。\\
\textbf{主要意义：}用于识别库存可行性与工作量均衡之间的关系；正式结论仍应回到分类缺口表。\\
\textbf{正文建议：}备选

\subsection*{图 23：K=3 总资源缺口—运输工作量均衡关系}
\textbf{文件：}\texttt{fig\_q4\_23\_K3\_shortage\_balance\_scatter.png}；矢量版：\texttt{fig\_q4\_23\_K3\_shortage\_balance\_scatter.pdf}\\
\textbf{图中内容：}按资源类别分别计算缺口后求和，仅用于散点横轴，不把不同资源成本等价化。\\
\textbf{主要意义：}用于识别库存可行性与工作量均衡之间的关系；正式结论仍应回到分类缺口表。\\
\textbf{正文建议：}备选

\subsection*{图 24：K=2 中继副本数量—新增能耗关系}
\textbf{文件：}\texttt{fig\_q4\_24\_K2\_copy\_energy\_scatter.png}；矢量版：\texttt{fig\_q4\_24\_K2\_copy\_energy\_scatter.pdf}\\
\textbf{图中内容：}每个点对应一个分区，新增能耗按被复制原中继架次的实际能耗累加。\\
\textbf{主要意义：}相同复制次数仍可能对应不同新增能耗，说明应使用实际架次能耗而不是统一副本惩罚系数。\\
\textbf{正文建议：}备选

\subsection*{图 25：K=3 中继副本数量—新增能耗关系}
\textbf{文件：}\texttt{fig\_q4\_25\_K3\_copy\_energy\_scatter.png}；矢量版：\texttt{fig\_q4\_25\_K3\_copy\_energy\_scatter.pdf}\\
\textbf{图中内容：}每个点对应一个分区，新增能耗按被复制原中继架次的实际能耗累加。\\
\textbf{主要意义：}相同复制次数仍可能对应不同新增能耗，说明应使用实际架次能耗而不是统一副本惩罚系数。\\
\textbf{正文建议：}备选

\subsection*{图 26：K=2主方案资源需求与库存对比}
\textbf{文件：}\texttt{fig\_q4\_26\_K2\_resource\_vs\_stock.png}；矢量版：\texttt{fig\_q4\_26\_K2\_resource\_vs\_stock.pdf}\\
\textbf{图中内容：}对8类资源分别比较任务组独立执行后的总需求与题目现有库存。\\
\textbf{主要意义：}直接识别具体资源缺口，避免用不同类型资源之间的无依据加权成本相互抵消。\\
\textbf{正文建议：}正文强烈推荐

\subsection*{图 27：K=3主方案资源需求与库存对比}
\textbf{文件：}\texttt{fig\_q4\_27\_K3\_resource\_vs\_stock.png}；矢量版：\texttt{fig\_q4\_27\_K3\_resource\_vs\_stock.pdf}\\
\textbf{图中内容：}对8类资源分别比较任务组独立执行后的总需求与题目现有库存。\\
\textbf{主要意义：}直接识别具体资源缺口，避免用不同类型资源之间的无依据加权成本相互抵消。\\
\textbf{正文建议：}正文强烈推荐

\subsection*{图 28：K=2各任务组独立资源配置矩阵}
\textbf{文件：}\texttt{fig\_q4\_28\_K2\_group\_resource\_heatmap.png}；矢量版：\texttt{fig\_q4\_28\_K2\_group\_resource\_heatmap.pdf}\\
\textbf{图中内容：}按任务组展示A/B/C运输机、电池、中继机和中继能源组件的独立最少配置。\\
\textbf{主要意义：}用于解释提交模板中每一行资源数量来自何处，并观察不同组之间的资源压力差异。\\
\textbf{正文建议：}正文推荐

\subsection*{图 29：K=3各任务组独立资源配置矩阵}
\textbf{文件：}\texttt{fig\_q4\_29\_K3\_group\_resource\_heatmap.png}；矢量版：\texttt{fig\_q4\_29\_K3\_group\_resource\_heatmap.pdf}\\
\textbf{图中内容：}按任务组展示A/B/C运输机、电池、中继机和中继能源组件的独立最少配置。\\
\textbf{主要意义：}用于解释提交模板中每一行资源数量来自何处，并观察不同组之间的资源压力差异。\\
\textbf{正文建议：}正文推荐

\subsection*{图 30：K=2主方案组间工作量}
\textbf{文件：}\texttt{fig\_q4\_30\_K2\_group\_workload.png}；矢量版：\texttt{fig\_q4\_30\_K2\_group\_workload.pdf}\\
\textbf{图中内容：}运输任务和中继副本的累计作业时间采用堆叠柱展示；运输任务部分不受中继复制口径影响。\\
\textbf{主要意义：}运输工作量CV=0.111，可用于评价分区的工作量均衡性，同时避免只按服务区数量判断均衡。\\
\textbf{正文建议：}正文强烈推荐

\subsection*{图 31：K=3主方案组间工作量}
\textbf{文件：}\texttt{fig\_q4\_31\_K3\_group\_workload.png}；矢量版：\texttt{fig\_q4\_31\_K3\_group\_workload.pdf}\\
\textbf{图中内容：}运输任务和中继副本的累计作业时间采用堆叠柱展示；运输任务部分不受中继复制口径影响。\\
\textbf{主要意义：}运输工作量CV=0.448，可用于评价分区的工作量均衡性，同时避免只按服务区数量判断均衡。\\
\textbf{正文建议：}正文强烈推荐

\subsection*{图 32：三种资源共享口径的配置规模比较}
\textbf{文件：}\texttt{fig\_q4\_32\_resource\_redundancy\_compare.png}；矢量版：\texttt{fig\_q4\_32\_resource\_redundancy\_compare.pdf}\\
\textbf{图中内容：}将问题三固定任务统一共享时的理论最少资源与2组、3组主方案独立配置需求并列。\\
\textbf{主要意义：}直接展示取消跨组资源复用造成的配置冗余，并可看出分组越细通信资源冗余越明显。\\
\textbf{正文建议：}正文强烈推荐

\subsection*{图 33：问题三资源编号与最少共享需求对比}
\textbf{文件：}\texttt{fig\_q4\_33\_original\_ids\_vs\_minimum.png}；矢量版：\texttt{fig\_q4\_33\_original\_ids\_vs\_minimum.pdf}\\
\textbf{图中内容：}在保持问题三所有任务时刻不变的情况下，重新做区间划分，比较原方案实际使用的资源编号数与理论最少数量。\\
\textbf{主要意义：}说明“出现过多少个编号”并不等于“最少需要多少资源”；B型电池和中继能源组件存在进一步复用空间。\\
\textbf{正文建议：}正文推荐

\subsection*{图 34：K=2主方案服务区归属条带图}
\textbf{文件：}\texttt{fig\_q4\_34\_K2\_service\_membership.png}；矢量版：\texttt{fig\_q4\_34\_K2\_service\_membership.pdf}\\
\textbf{图中内容：}按服务区编号顺序，用颜色标识15个服务区的任务组归属。\\
\textbf{主要意义：}作为任务分区的紧凑总览图，适合与地图配合使用；即使没有DEM也能清楚复核提交结果。\\
\textbf{正文建议：}正文备选

\subsection*{图 35：K=3主方案服务区归属条带图}
\textbf{文件：}\texttt{fig\_q4\_35\_K3\_service\_membership.png}；矢量版：\texttt{fig\_q4\_35\_K3\_service\_membership.pdf}\\
\textbf{图中内容：}按服务区编号顺序，用颜色标识15个服务区的任务组归属。\\
\textbf{主要意义：}作为任务分区的紧凑总览图，适合与地图配合使用；即使没有DEM也能清楚复核提交结果。\\
\textbf{正文建议：}正文备选

\subsection*{图 36：K=2中继任务按组复制矩阵}
\textbf{文件：}\texttt{fig\_q4\_36\_K2\_relay\_replication.png}；矢量版：\texttt{fig\_q4\_36\_K2\_relay\_replication.pdf}\\
\textbf{图中内容：}行表示问题三的原中继架次，列表示任务组；单元格非零表示该组需要一个保持原参数不变的中继任务副本。\\
\textbf{主要意义：}可逐项解释总复制架次和新增中继能耗，并验证同一原中继在一个组内最多复制一次。\\
\textbf{正文建议：}正文推荐

\subsection*{图 37：K=3中继任务按组复制矩阵}
\textbf{文件：}\texttt{fig\_q4\_37\_K3\_relay\_replication.png}；矢量版：\texttt{fig\_q4\_37\_K3\_relay\_replication.pdf}\\
\textbf{图中内容：}行表示问题三的原中继架次，列表示任务组；单元格非零表示该组需要一个保持原参数不变的中继任务副本。\\
\textbf{主要意义：}可逐项解释总复制架次和新增中继能耗，并验证同一原中继在一个组内最多复制一次。\\
\textbf{正文建议：}正文推荐

\subsection*{图 38：主方案资源冗余来源}
\textbf{文件：}\texttt{fig\_q4\_38\_redundancy\_source.png}；矢量版：\texttt{fig\_q4\_38\_redundancy\_source.pdf}\\
\textbf{图中内容：}将相对统一共享最少需求增加的资源，按运输侧和通信侧两大类汇总展示。\\
\textbf{主要意义：}主方案优先维持运输侧零增配，因此新增配置主要集中在共享中继被拆分后的通信侧。\\
\textbf{正文建议：}正文推荐

\subsection*{图 39：K=2 各跨组中继任务的新增复制能耗}
\textbf{文件：}\texttt{fig\_q4\_39\_K2\_copy\_energy.png}；矢量版：\texttt{fig\_q4\_39\_K2\_copy\_energy.pdf}\\
\textbf{图中内容：}仅展示跨组超边；柱高等于(涉及组数-1)×原中继架次能耗。\\
\textbf{主要意义：}全部柱高之和为2.471 kWh，是完整复制假设下可直接计算的中继能耗增量。\\
\textbf{正文建议：}正文推荐

\subsection*{图 40：K=3 各跨组中继任务的新增复制能耗}
\textbf{文件：}\texttt{fig\_q4\_40\_K3\_copy\_energy.png}；矢量版：\texttt{fig\_q4\_40\_K3\_copy\_energy.pdf}\\
\textbf{图中内容：}仅展示跨组超边；柱高等于(涉及组数-1)×原中继架次能耗。\\
\textbf{主要意义：}全部柱高之和为3.918 kWh，是完整复制假设下可直接计算的中继能耗增量。\\
\textbf{正文建议：}正文推荐

\subsection*{图 41：K=2固定运输任务按组甘特图}
\textbf{文件：}\texttt{fig\_q4\_41\_K2\_transport\_gantt.png}；矢量版：\texttt{fig\_q4\_41\_K2\_transport\_gantt.pdf}\\
\textbf{图中内容：}不改变问题三的任何运输时刻，只按最终任务组重新排列展示25个运输架次。\\
\textbf{主要意义：}用于复核分区没有改变原运输调度，并观察各组的运输高峰与并发结构。\\
\textbf{正文建议：}附录/备选

\subsection*{图 42：K=3固定运输任务按组甘特图}
\textbf{文件：}\texttt{fig\_q4\_42\_K3\_transport\_gantt.png}；矢量版：\texttt{fig\_q4\_42\_K3\_transport\_gantt.pdf}\\
\textbf{图中内容：}不改变问题三的任何运输时刻，只按最终任务组重新排列展示25个运输架次。\\
\textbf{主要意义：}用于复核分区没有改变原运输调度，并观察各组的运输高峰与并发结构。\\
\textbf{正文建议：}附录/备选

\subsection*{图 43：K=2中继副本甘特图}
\textbf{文件：}\texttt{fig\_q4\_43\_K2\_relay\_gantt.png}；矢量版：\texttt{fig\_q4\_43\_K2\_relay\_gantt.pdf}\\
\textbf{图中内容：}每组列出其所需的完整中继任务副本，中继无人机占用时段延续到返回O01后300秒。\\
\textbf{主要意义：}直观说明为什么“新增中继架次数”不等于“新增中继无人机数量”：不同时段副本可由同一组内实体中继机顺序执行。\\
\textbf{正文建议：}正文推荐

\subsection*{图 44：K=3中继副本甘特图}
\textbf{文件：}\texttt{fig\_q4\_44\_K3\_relay\_gantt.png}；矢量版：\texttt{fig\_q4\_44\_K3\_relay\_gantt.pdf}\\
\textbf{图中内容：}每组列出其所需的完整中继任务副本，中继无人机占用时段延续到返回O01后300秒。\\
\textbf{主要意义：}直观说明为什么“新增中继架次数”不等于“新增中继无人机数量”：不同时段副本可由同一组内实体中继机顺序执行。\\
\textbf{正文建议：}正文推荐

\subsection*{图 45：K=2均衡阈值敏感性}
\textbf{文件：}\texttt{fig\_q4\_45\_K2\_balance\_sensitivity.png}；矢量版：\texttt{fig\_q4\_45\_K2\_balance\_sensitivity.pdf}\\
\textbf{图中内容：}逐步收紧运输工作量CV上限，重新统计满足该均衡要求的分区中最小中继无人机需求。\\
\textbf{主要意义：}用于说明过强的均衡要求可能迫使更多共享中继被切开，从而提高通信资源配置下界。\\
\textbf{正文建议：}正文备选

\subsection*{图 46：K=3均衡阈值敏感性}
\textbf{文件：}\texttt{fig\_q4\_46\_K3\_balance\_sensitivity.png}；矢量版：\texttt{fig\_q4\_46\_K3\_balance\_sensitivity.pdf}\\
\textbf{图中内容：}逐步收紧运输工作量CV上限，重新统计满足该均衡要求的分区中最小中继无人机需求。\\
\textbf{主要意义：}用于说明过强的均衡要求可能迫使更多共享中继被切开，从而提高通信资源配置下界。\\
\textbf{正文建议：}正文备选

\subsection*{图 47：K=2前25候选方案对比}
\textbf{文件：}\texttt{fig\_q4\_47\_K2\_top\_candidates.png}；矢量版：\texttt{fig\_q4\_47\_K2\_top\_candidates.pdf}\\
\textbf{图中内容：}按主筛选规则取前25个候选，纵轴为中继无人机需求，颜色表示运输工作量CV。\\
\textbf{主要意义：}展示最终方案附近仍存在的资源—均衡权衡，避免只报告一个方案而看不到邻近候选结构。\\
\textbf{正文建议：}附录/备选

\subsection*{图 48：K=3前25候选方案对比}
\textbf{文件：}\texttt{fig\_q4\_48\_K3\_top\_candidates.png}；矢量版：\texttt{fig\_q4\_48\_K3\_top\_candidates.pdf}\\
\textbf{图中内容：}按主筛选规则取前25个候选，纵轴为中继无人机需求，颜色表示运输工作量CV。\\
\textbf{主要意义：}展示最终方案附近仍存在的资源—均衡权衡，避免只报告一个方案而看不到邻近候选结构。\\
\textbf{正文建议：}附录/备选

\subsection*{图 49：K=2 主方案库存利用率}
\textbf{文件：}\texttt{fig\_q4\_49\_K2\_inventory\_utilization.png}；矢量版：\texttt{fig\_q4\_49\_K2\_inventory\_utilization.pdf}\\
\textbf{图中内容：}按8类资源计算主方案总需求与现有库存的比例。\\
\textbf{主要意义：}超过100\%的资源就是必须增配的类别；同时可观察未形成缺口的资源冗余程度。\\
\textbf{正文建议：}正文备选

\subsection*{图 50：K=3 主方案库存利用率}
\textbf{文件：}\texttt{fig\_q4\_50\_K3\_inventory\_utilization.png}；矢量版：\texttt{fig\_q4\_50\_K3\_inventory\_utilization.pdf}\\
\textbf{图中内容：}按8类资源计算主方案总需求与现有库存的比例。\\
\textbf{主要意义：}超过100\%的资源就是必须增配的类别；同时可观察未形成缺口的资源冗余程度。\\
\textbf{正文建议：}正文备选

\subsection*{图 51：K=2各组任务规模比较}
\textbf{文件：}\texttt{fig\_q4\_51\_K2\_group\_task\_scale.png}；矢量版：\texttt{fig\_q4\_51\_K2\_group\_task\_scale.pdf}\\
\textbf{图中内容：}同时比较各组服务区数、固定运输架次数和货箱数。\\
\textbf{主要意义：}强调“服务区数量均衡”不等于“实际工作量均衡”，为采用累计任务时间CV提供解释。\\
\textbf{正文建议：}正文备选

\subsection*{图 52：K=3各组任务规模比较}
\textbf{文件：}\texttt{fig\_q4\_52\_K3\_group\_task\_scale.png}；矢量版：\texttt{fig\_q4\_52\_K3\_group\_task\_scale.pdf}\\
\textbf{图中内容：}同时比较各组服务区数、固定运输架次数和货箱数。\\
\textbf{主要意义：}强调“服务区数量均衡”不等于“实际工作量均衡”，为采用累计任务时间CV提供解释。\\
\textbf{正文建议：}正文备选

\subsection*{图 53：K=2主方案空间分区图}
\textbf{文件：}\texttt{fig\_q4\_53\_K2\_geo\_partition.png}；矢量版：\texttt{fig\_q4\_53\_K2\_geo\_partition.pdf}\\
\textbf{图中内容：}按真实服务区经纬度绘制最终任务组；若可读取DEM则以地形灰度作为背景。\\
\textbf{主要意义：}用于判断任务组的空间分布和紧凑程度。地理距离只作为辅助解释，不覆盖固定运输与中继关联约束。\\
\textbf{正文建议：}正文强烈推荐

\subsection*{图 54：K=3主方案空间分区图}
\textbf{文件：}\texttt{fig\_q4\_54\_K3\_geo\_partition.png}；矢量版：\texttt{fig\_q4\_54\_K3\_geo\_partition.pdf}\\
\textbf{图中内容：}按真实服务区经纬度绘制最终任务组；若可读取DEM则以地形灰度作为背景。\\
\textbf{主要意义：}用于判断任务组的空间分布和紧凑程度。地理距离只作为辅助解释，不覆盖固定运输与中继关联约束。\\
\textbf{正文建议：}正文强烈推荐

\subsection*{图 55：K=2任务组与中继悬停点}
\textbf{文件：}\texttt{fig\_q4\_55\_K2\_relay\_positions.png}；矢量版：\texttt{fig\_q4\_55\_K2\_relay\_positions.pdf}\\
\textbf{图中内容：}同一颜色同时标记任务组服务区和该组需要复制的原中继悬停位置。\\
\textbf{主要意义：}揭示分区后通信保障的空间复制范围，并可解释某些组虽然服务区较少仍需要较多中继资源。\\
\textbf{正文建议：}正文备选

\subsection*{图 56：K=3任务组与中继悬停点}
\textbf{文件：}\texttt{fig\_q4\_56\_K3\_relay\_positions.png}；矢量版：\texttt{fig\_q4\_56\_K3\_relay\_positions.pdf}\\
\textbf{图中内容：}同一颜色同时标记任务组服务区和该组需要复制的原中继悬停位置。\\
\textbf{主要意义：}揭示分区后通信保障的空间复制范围，并可解释某些组虽然服务区较少仍需要较多中继资源。\\
\textbf{正文建议：}正文备选

\subsection*{图 57：K=2固定运输路线与空间分区}
\textbf{文件：}\texttt{fig\_q4\_57\_K2\_geo\_routes.png}；矢量版：\texttt{fig\_q4\_57\_K2\_geo\_routes.pdf}\\
\textbf{图中内容：}在30 m DEM背景上按任务组颜色叠加问题三已经固定的25个运输架次路线；任何路线均未因问题四重新优化。\\
\textbf{主要意义：}直观验证同一多点运输架次涉及的服务区位于同一组，同时展示空间分区与既有运输网络的关系。\\
\textbf{正文建议：}正文强烈推荐

\subsection*{图 58：K=3固定运输路线与空间分区}
\textbf{文件：}\texttt{fig\_q4\_58\_K3\_geo\_routes.png}；矢量版：\texttt{fig\_q4\_58\_K3\_geo\_routes.pdf}\\
\textbf{图中内容：}在30 m DEM背景上按任务组颜色叠加问题三已经固定的25个运输架次路线；任何路线均未因问题四重新优化。\\
\textbf{主要意义：}直观验证同一多点运输架次涉及的服务区位于同一组，同时展示空间分区与既有运输网络的关系。\\
\textbf{正文建议：}正文强烈推荐

\subsection*{图 59：K=2运输分区与中继复制空间综合图}
\textbf{文件：}\texttt{fig\_q4\_59\_K2\_geo\_routes\_relays.png}；矢量版：\texttt{fig\_q4\_59\_K2\_geo\_routes\_relays.pdf}\\
\textbf{图中内容：}在DEM背景上同时叠加固定运输路线、任务组服务区和各组所需的原中继悬停位置；空心三角表示中继任务按组完整复制。\\
\textbf{主要意义：}强调复制只增加执行资源，不移动中继位置、不改变海拔和服务时段，因此问题三链路几何与链路预算可直接继承。\\
\textbf{正文建议：}正文强烈推荐

\subsection*{图 60：K=3运输分区与中继复制空间综合图}
\textbf{文件：}\texttt{fig\_q4\_60\_K3\_geo\_routes\_relays.png}；矢量版：\texttt{fig\_q4\_60\_K3\_geo\_routes\_relays.pdf}\\
\textbf{图中内容：}在DEM背景上同时叠加固定运输路线、任务组服务区和各组所需的原中继悬停位置；空心三角表示中继任务按组完整复制。\\
\textbf{主要意义：}强调复制只增加执行资源，不移动中继位置、不改变海拔和服务时段，因此问题三链路几何与链路预算可直接继承。\\
\textbf{正文建议：}正文强烈推荐

\subsection*{图 61：K=2任务组1空间执行细节}
\textbf{文件：}\texttt{fig\_q4\_61\_K2\_group1\_geo\_detail.png}；矢量版：\texttt{fig\_q4\_61\_K2\_group1\_geo\_detail.pdf}\\
\textbf{图中内容：}单独放大任务组1，显示其服务区、固定运输路线以及需要保留/复制的原中继悬停位置。\\
\textbf{主要意义：}用于逐组解释资源配置来源，尤其适合附录核查某组为何需要特定数量的中继无人机和能源组件。\\
\textbf{正文建议：}附录/备选

\subsection*{图 62：K=2任务组2空间执行细节}
\textbf{文件：}\texttt{fig\_q4\_62\_K2\_group2\_geo\_detail.png}；矢量版：\texttt{fig\_q4\_62\_K2\_group2\_geo\_detail.pdf}\\
\textbf{图中内容：}单独放大任务组2，显示其服务区、固定运输路线以及需要保留/复制的原中继悬停位置。\\
\textbf{主要意义：}用于逐组解释资源配置来源，尤其适合附录核查某组为何需要特定数量的中继无人机和能源组件。\\
\textbf{正文建议：}附录/备选

\subsection*{图 63：K=3任务组1空间执行细节}
\textbf{文件：}\texttt{fig\_q4\_63\_K3\_group1\_geo\_detail.png}；矢量版：\texttt{fig\_q4\_63\_K3\_group1\_geo\_detail.pdf}\\
\textbf{图中内容：}单独放大任务组1，显示其服务区、固定运输路线以及需要保留/复制的原中继悬停位置。\\
\textbf{主要意义：}用于逐组解释资源配置来源，尤其适合附录核查某组为何需要特定数量的中继无人机和能源组件。\\
\textbf{正文建议：}附录/备选

\subsection*{图 64：K=3任务组2空间执行细节}
\textbf{文件：}\texttt{fig\_q4\_64\_K3\_group2\_geo\_detail.png}；矢量版：\texttt{fig\_q4\_64\_K3\_group2\_geo\_detail.pdf}\\
\textbf{图中内容：}单独放大任务组2，显示其服务区、固定运输路线以及需要保留/复制的原中继悬停位置。\\
\textbf{主要意义：}用于逐组解释资源配置来源，尤其适合附录核查某组为何需要特定数量的中继无人机和能源组件。\\
\textbf{正文建议：}附录/备选

\subsection*{图 65：K=3任务组3空间执行细节}
\textbf{文件：}\texttt{fig\_q4\_65\_K3\_group3\_geo\_detail.png}；矢量版：\texttt{fig\_q4\_65\_K3\_group3\_geo\_detail.pdf}\\
\textbf{图中内容：}单独放大任务组3，显示其服务区、固定运输路线以及需要保留/复制的原中继悬停位置。\\
\textbf{主要意义：}用于逐组解释资源配置来源，尤其适合附录核查某组为何需要特定数量的中继无人机和能源组件。\\
\textbf{正文建议：}附录/备选
